\clearpage\section{MOS 電晶體電路圖}
\begingroup\small
來源 netlist：\texttt{01\_counters.sp}. 同步與漣波計數器的 MOS 單元與完整階層。

相同節點標籤代表電氣相連。下圖為可編輯的 CircuitikZ 原始碼；階層方塊均對應本附錄所畫出的 MOS 單元。

除另有標示外，所有 NMOS 體端接 \texttt{vn}=0，所有 PMOS 體端接 \texttt{vp}=\(V_{DD}\)；為保持可讀性，省略體端導線。一般反相器及動態／時脈單元為 \(W_n/L_n=0.6/0.25\,\mu\mathrm m\)、\(W_p/L_p=1.8/0.25\,\mu\mathrm m\)；TG 兩顆管皆為 \(W/L=0.6/0.25\,\mu\mathrm m\)。

\noindent\begin{minipage}[t]{.48\linewidth}\centering
\begin{circuitikz}[scale=0.92, transform shape, font=\small]
\node[pmos] (MP) at (0,0.0) {};
\draw (MP.G) -- ++(-.65,0) node[left] {\texttt{a}};
\node[right=3pt] at (MP.right) {\texttt{MP}};
\node[nmos] (MN) at (0,-1.45) {};
\draw (MN.G) -- ++(-.65,0) node[left] {\texttt{a}};
\node[right=3pt] at (MN.right) {\texttt{MN}};
\draw (MP.S) -- ++(0,.55) coordinate (net0) node[above] {\texttt{vp}};
\draw (MP.D) -- (MN.D) coordinate[midway] (net1);
\draw (net1) node[circ] {} -- ++(1.0,0) node[right] {\texttt{y}};
\draw (MN.S) -- ++(0,-.55) coordinate (net2) node[below] {\texttt{vn}};

\end{circuitikz}
\par\small \texttt{inv}: 2T
\end{minipage}\hfill\begin{minipage}[t]{.48\linewidth}\centering
\begin{circuitikz}[scale=0.88, transform shape, font=\small]
\node[nmos] (MN) at (0,0) {};
\draw (MN.G) -- ++(-.65,0) node[left] {\texttt{en}};
\node[right=3pt] at (MN.right) {\texttt{MN}};
\node[pmos,yscale=-1] (MP) at (3.4,0) {};
\draw (MP.G) -- ++(-.65,0) node[left] {\texttt{eb}};
\node[right=3pt] at (MP.right) {\texttt{MP}};

\draw (MN.D) -- ++(0,.8) coordinate (nt) -- (nt -| MP.D) -- (MP.D);
\draw (MN.S) -- ++(0,-.8) coordinate (nb) -- (nb -| MP.S) -- (MP.S);
\node[above] at ($(nt)!.5!(nt -| MP.D)$) {\texttt{b}};
\node[below] at ($(nb)!.5!(nb -| MP.S)$) {\texttt{a}};

\end{circuitikz}
\par\small \texttt{tg}: 2T
\end{minipage}\par\medskip
\noindent\begin{minipage}[t]{.48\linewidth}\centering
\begin{circuitikz}[scale=0.82, transform shape, font=\small]
\node[pmos] (MP1) at (0,0) {};
\draw (MP1.G) -- ++(-.65,0) node[left] {\texttt{a}};
\node[right=3pt] at (MP1.right) {\texttt{MP1}};
\node[pmos] (MP2) at (3.8,0) {};
\draw (MP2.G) -- ++(-.65,0) node[left] {\texttt{b}};
\node[right=3pt] at (MP2.right) {\texttt{MP2}};
\node[nmos] (MN1) at (0,-2) {};
\draw (MN1.G) -- ++(-.65,0) node[left] {\texttt{a}};
\node[right=3pt] at (MN1.right) {\texttt{MN1}};
\node[nmos] (MN2) at (0,-3.5) {};
\draw (MN2.G) -- ++(-.65,0) node[left] {\texttt{b}};
\node[right=3pt] at (MN2.right) {\texttt{MN2}};

\draw (MP1.S) -- ++(0,.55) coordinate (v) -- (v -| MP2.S) -- (MP2.S);
\node[above] at (v) {\texttt{vp}};
\draw (MP1.D) -- (MN1.D) coordinate[midway] (y);
\draw (MP2.D) -- (MP2.D |- y) -- (y) node[circ] {};
\node[right] at (MP2.D |- y) {\texttt{y}};
\draw (MN1.S) -- (MN2.D) node[midway,right] {\texttt{x}};
\draw (MN2.S) -- ++(0,-.55) node[below] {\texttt{vn}};

\end{circuitikz}
\par\small \texttt{nand2}: 4T; NMOS \(1.2/0.25\), PMOS \(1.8/0.25\).
\end{minipage}\hfill\begin{minipage}[t]{.48\linewidth}\centering
\begin{circuitikz}[scale=0.82, transform shape, font=\small]
\node[pmos] (MP1) at (0,0) {};
\draw (MP1.G) -- ++(-.65,0) node[left] {\texttt{rst}};
\node[right=3pt] at (MP1.right) {\texttt{MP1}};
\node[pmos] (MP2) at (0,-1.5) {};
\draw (MP2.G) -- ++(-.65,0) node[left] {\texttt{z}};
\node[right=3pt] at (MP2.right) {\texttt{MP2}};
\node[nmos] (MN1) at (0,-3.5) {};
\draw (MN1.G) -- ++(-.65,0) node[left] {\texttt{z}};
\node[right=3pt] at (MN1.right) {\texttt{MN1}};
\node[nmos] (MN2) at (3.8,-3.5) {};
\draw (MN2.G) -- ++(-.65,0) node[left] {\texttt{rst}};
\node[right=3pt] at (MN2.right) {\texttt{MN2}};

\draw (MP1.S) -- ++(0,.55) node[above] {\texttt{vp}};
\draw (MP1.D) -- (MP2.S) node[midway,right] {\texttt{pr}};
\draw (MP2.D) -- (MN1.D) coordinate[midway] (q);
\draw (q) node[circ] {} -- (q -| MN2.D) -- (MN2.D);
\node[right] at (q -| MN2.D) {\texttt{q}};
\draw (MN1.S) -- ++(0,-.55) coordinate (g) -- (g -| MN2.S) -- (MN2.S);
\node[below] at (g) {\texttt{vn}};

\end{circuitikz}
\par\small 計數器從鎖存器 NOR：4T；NMOS \(0.6/0.25\)，PMOS \(3.6/0.25\)。
\end{minipage}\par\medskip
\clearpage\subsection{具有重置功能的 22 顆電晶體儲存單元}
\begin{center}
\begin{circuitikz}[scale=0.94, transform shape, font=\small]
\node[draw,rounded corners=1pt,minimum width=2.6cm,minimum height=1cm,align=center] (b0) at (0.0,0) {XT1: TG\\en=cb, eb=ck};
\node[draw,rounded corners=1pt,minimum width=2.6cm,minimum height=1cm,align=center] (b1) at (3.25,0) {XI1: NAND\\inputs x, rb};
\draw[-{Latex}] (b0.east) -- node[above,font=\scriptsize] {x} (b1.west);
\node[draw,rounded corners=1pt,minimum width=2.6cm,minimum height=1cm,align=center] (b2) at (6.5,0) {XT3: TG\\en=ck, eb=cb};
\draw[-{Latex}] (b1.east) -- node[above,font=\scriptsize] {y} (b2.west);
\node[draw,rounded corners=1pt,minimum width=2.6cm,minimum height=1cm,align=center] (b3) at (9.75,0) {4T NOR\\inputs z, rst};
\draw[-{Latex}] (b2.east) -- node[above,font=\scriptsize] {z} (b3.west);

\end{circuitikz}
\end{center}
\begin{center}\small\begin{tabular}{p{.24\linewidth}p{.68\linewidth}}\toprule
部分 & 確切接線\\\midrule
XR & INV(\texttt{rst}\(\to\)\texttt{rb}).\\
Master feedback & XI2: INV(\texttt{y}\(\to\)\texttt{fb1}); XT2: TG(\texttt{fb1}\(\leftrightarrow\)\texttt{x}, en=ck, eb=cb).\\
Slave feedback & XI4: INV(\texttt{q}\(\to\)\texttt{fb2}); XT4: TG(\texttt{fb2}\(\leftrightarrow\)\texttt{z}, en=cb, eb=ck).\\
Input / output & XT1 input=\texttt{d}; NOR output=\texttt{q}.\newline \texttt{counterff(d,ck,cb,rst,q,x,y,z,vp,vn)}.\\
\bottomrule\end{tabular}\end{center}
重置同時作用於兩個儲存級：\(y=\overline{x\,rb}\)、\(q=\overline{z+rst}\)。此單元包含四個 TG（8T）、三個反相器（6T）、一個 NAND（4T）與一個 NOR（4T）。
\begin{center}\small\begin{tabular}{p{.24\linewidth}p{.68\linewidth}}\toprule
複合單元 & 完整 MOS 單元展開\\\midrule
\texttt{and2(a,b,y)} & NAND(a,b,z); INV(z,y).\\
\texttt{xor2(a,b,y)} & NAND(a,b,n1); NAND(a,n1,n2); NAND(b,n1,n3); NAND(n2,n3,y).\\
\texttt{syncgate(ck,en,y)} & NAND(ck,n2,n1); NAND(n1,en,n2); NAND(n2,n4,n3); NAND(n3,ck,n4); INV(n4,y).\\
\bottomrule\end{tabular}\end{center}
上表 NAND 與 INV 參數順序為輸入後接輸出；所有單元共用所屬支路的電源與接地。
\clearpage\subsection{計數器接線}
\begin{center}
\begin{circuitikz}[scale=0.94, transform shape, font=\small]
\node[draw,rounded corners=1pt,minimum width=2.6cm,minimum height=1cm,align=center] (b0) at (0.0,0) {XJ1: counterff\\D=j2b, Q=j1};
\node[draw,rounded corners=1pt,minimum width=2.6cm,minimum height=1cm,align=center] (b1) at (3.25,0) {XJ2: counterff\\D=j1, Q=j2};
\draw[-{Latex}] (b0.east) -- node[above,font=\scriptsize] {j1} (b1.west);
\node[draw,rounded corners=1pt,minimum width=2.6cm,minimum height=1cm,align=center] (b2) at (6.5,0) {XJ3: counterff\\D=j2, Q=b1};
\draw[-{Latex}] (b1.east) -- node[above,font=\scriptsize] {j2} (b2.west);

\end{circuitikz}
\end{center}
三個 Johnson 單元共用 \texttt{ck=clk, cb=cb}。\texttt{XJC}：INV(j2,j2b)；\texttt{XJB}：XOR(b1,j2b,b0)。因此 \(D_1=\overline{J_2}\)、\(D_2=J_1\)、\(D_3=J_2\)、\(b_1=J_3\)、\(b_0=J_3\oplus\overline{J_2}\)，保留講義特殊的回授接法。重置後 \(J_1J_2J_3=000\) 解碼為 \(b_1b_0=01\)，穩態低位循環為 \(1,0,3,2\)。
\begin{center}\small\begin{tabular}{p{.24\linewidth}p{.68\linewidth}}\toprule
網路 & 完整階層\\\midrule
Clock enables & XE1: AND(b0,b1,en1); XE2: AND(b2,b3,en23); XE3: AND(en23,b4,en234); XE4: AND(en1,en234,en2).\\
Gated clocks & XG1: syncgate(clk,en1,ck2); XG2: syncgate(clk,en2,ck5). XCB2: INV(ck2,cb2); XCB5: INV(ck5,cb5). XCS: INV(clk,cb).\\
Bits 2--4 & For \(i=2,3,4\): \(d_i=b_i\oplus c_{i-1}\), \(c_i=b_i c_{i-1}\), \(c_1=1\). Each counterff uses \texttt{ck2,cb2}, output \(b_i\).\\
Bits 5--7 & For \(i=5,6,7\): \(d_i=b_i\oplus c_{i-1}\), \(c_i=b_i c_{i-1}\), fresh \(c_4=1\). Each counterff uses \texttt{ck5,cb5}, output \(b_i\).\\
Ripple bits 0--3 & XAR\(i\): INV(\(a_i,ab_i\)); XA\(i\): counterff with \(D=ab_i,Q=a_i\). XA0 uses \texttt{clk,cba}, XCA: INV(clk,cba). XA\(i\), \(i>0\), uses \(ck=ab_{i-1},cb=a_{i-1}\).\\
\bottomrule\end{tabular}\end{center}
所有儲存單元共用 \texttt{rst}。同步支路電源 \texttt{vs} 與漣波支路電源 \texttt{va} 為獨立 1 V 電源，以便分別量測電流。每個位元輸出皆接 5 fF 至 0；上述階層展開後，同步支路為 414 顆 MOS，漣波支路為 98 顆。
\endgroup
